\section{Постановка проблемы}

\subsection{Основные обозначения}
В дальнейшем будут использоваться следующие основные обозначения. Пусть $X = \left( x_i^j \right) \in \mathbb{R}^{n \times d}$ обозначает матрицу наблюдений размером $n \times d$, где каждый $d$-мерный объект $x_i = (x_i^1, x_i^2, \dots, x_i^d) \in \mathbb{R}^{1 \times d}$ представляет собой строковый вектор признаков для каждого наблюдения $i=1,2,\dots,n$. Вектор $x_j = (x_1^j, x_2^j, \dots, x_n^j)^\top \in \mathbb{R}^{n \times 1}$ означает столбец значений признака $j$ для всех наблюдений $j=1,2,\dots,d$. Вектор $y = (y_1, y_2, \dots, y_n)^\top \in \mathbb{R}^{n \times 1}$ содержит значения целевой функции, представляющей искомые кластеры. Число кластеров обозначено как $k$, а разбиение наблюдений на кластеры — как $S={S_1, S_2, \dots, S_k}$, где каждый кластер $S_l \subset {1,2, \dots, n }$ для $l=1,2,\dots,k$. Набор из $k$ центров кластеров представлен через $c = { c_1,c_2, \dots, c_k }$, где каждый центр $c_l = (c_l^1, c_l^2, \dots, c_l^d) \in \mathbb{R}^{1 \times d}$ для $l=1,2,\dots,k$. Также используется диагональная матрица размерности $d$, обозначенная как $I_d$.


\subsection{Формулирование задачи кластеризации без учителя}
Задача кластеризации относится к типу задач <<без учителя>> (Unsupervised learning). Задача кластеризации -- есть группировка объектов в наборы (группы), основываясь только на их сходстве между собой так, чтобы объекты внутри одной группы были больше похожи на представителей этой группы, чем другой.

Формальная запись задач имеет следующий вид:
Для каждого $d$-мерного объекта $x_i$ из множества объектов $X = {x_1, x_2, ..., x_n}$, необходимо поставить в соответствие один из $k$ кластеров $c= \{c_1, c_2, \dots, c_k\}$.

Для оценивания <<схожести>> объектов могут быть использованы разные метрики расстояния такие как Евклидово расстояние \cite{Lloyd1982kmeans}, Манхеттенское расстояние, \cite{kaufman2008kmedioids}, расстояние Минковского \cite{aggarwal2001surprising} и другие. А также различные способы оценки межкласторных расстояний, такие как Average Linkage \cite{Finkelman2021averagelinkage}, Complete Linkage \cite{Defays1977CLINK}, Single Linkage \cite{Sibson1973SLINK}.


Задача кластеризации на основе центроидов для расстояния Минковского с параметром \(p\) есть задача сложной невыпуклой оптимизации функционала качества (\ref{eq:kmeans_problem}), которая может быть решена с помощью эвристического алгоритма \kmeans{} для метрики Минковского, где необходимо оптимизировать сумму расстояний между объектами внутри кластеров.

\begin{equation}
    \begin{gathered}
        W_p(S;c) = \sum_{l=1}^k \sum_{x \in S_l} \sum_{j=1}^{d} |x^j - c_l^j|^p \to \min, \\
        x=(x^1, x^2, \dots, x^d) \in R^{1 \times d}, \\
        c_l=(c_l^1,c_l^2, \dots, c_l^d) \in R^{1 \times d}
    \end{gathered}
    \label{eq:kmeans_problem}
\end{equation}

\subsubsection{Выбор метода кластеризации}

В данной работе основой для алгоритма кластеризации с учителем служит реализация \kmeans{} с метрикой Минковского для пространств большой размерности, поскольку \kmeans{} -- легко интерпретируемый подход, основанный на оптимизации суммарного квадратичного отклонения внутри кластеров. Выбор \kmeans{} также обусловлен его основными преимуществами: простотой реализации (существует большое число реализаций и оптимизаций данного подхода), высокой скоростью работы (на небольших наборах данных алгоритм способен сходится за очень короткое время). Алгоритм \kmeans{} выполняет кластеризацию данных, оптимизируя суммарное квадратичное отклонение внутри кластеров. Его работа состоит из нескольких этапов:
\begin{enumerate}
    \item Инициализация.
    Фиксируется параметр $k$ -- число кластеров, которые необходимо найти, после чего центроиды кластеров выбираются случайным образом.
    \item Присвоение данных кластерам.
    Каждый объект данных присваивается кластеру, чей центроид на основе выбранной метрики расстояния окажется ближе всего.
    \item Обновление центроидов.
    Центроиды пересчитываются на основе объектов из соответствующих кластеров.
    \item Итеративное повторение шагов 1, 2 и 3.
\end{enumerate}


\subsubsection{Выбор метрики качества}
Для оценки качества предложенного алгоритма кластеризации были выбраны метрики $Accuracy$, $NMI$, $ARI$, $AMI$.

\paragraph{Rand Index ($RI$)} является метрикой, используемой для измерения сходства между двумя различными кластеризациями или разбиениями одного и того же набора данных. Основная идея $RI$ заключается в том, чтобы подсчитать долю пар объектов, которые находятся в одном и том же кластере в обоих разбиениях (кластеризациях).
Для того чтобы определить $RI$, вводятся следующие обозначения. Пусть $X$ -- набор, состоящий из $n$ объектов для кластеризации, в котором $a$ пар объектов попали в один кластер, а $b$ пар попали в разные кластеры. Тогда метрика Rand Index ($RI$) будет иметь вид:
\begin{equation}
 RI = \frac{a+b}{C_2^n}.
\end{equation}

\paragraph{Adjusted Rand Index ($ARI$)} -- скорректированная версия Rand Index, предназначенная для измерения сходства между двумя кластеризациями, путем корректировки метрики $RI$. $ARI$ может принимать значения в промежутке $[0; 1]$, где значение 1 соответствует идентичным кластеризациям, а значение 0 означает, что кластеризации случайны.
\begin{equation}
    ARI = \frac{RI - E[RI]}{max(RI)- E[RI]},
\end{equation}
где $E[RI]$ - оценка математического ожидания $RI$, а $max(RI)$ - максимальное значение Rand Index, полученное в результате кластеризаций.

\paragraph{Mutual Information ($MI$)} измеряет степени зависимости между двумя различными разбиениями одного и того же набора данных на кластеры. Обозначим за $U$ и $V$ два разбиения набора данных $X$, состоящего из $n$ объектов, что $|U_i|$ -- есть количество объектов в кластере $U_i$, а $|V_j|$ в кластере $V_j$. Обозначим энтропию за $H(U)$.
\begin{equation}
    MI(U, V) = \displaystyle\sum_{i=1}^{|U|} \displaystyle\sum_{j=1}^{|V|} \frac{|U_i \cap V_j|}{n} \log\frac{n|U_i \cap V_i|}{|U_i||V_i|},
\end{equation}


\paragraph{Normalized Mutual Information ($NMI$)} -- вариант $MI$, который нормализуется для того, чтобы значения можно было интерпретировать на отрезке от 0 до 1. Метрика используется для измерения сходства между двумя различными разбиениями одного и того же набора данных на кластеры.
\begin{equation}
    NMI(U, V) = \frac{MI(U, V)}{mean(H(U), H(V))}.
\end{equation}

\paragraph{Adjusted Mutual Information ($AMI$)} нормализуется относительно максимальной величины, которую может принять взаимная информация, что позволяет оценивать сходство между разбиениями на независимых от размера кластеров данных. Она принимает значения в интервале от 0 до 1, где значение 1 означает идентичные разбиения, а значение 0 указывает на полное отсутствие сходства. Данная метрика широко используется в оценке качества различных методов кластеризации или сегментации данных, особенно когда требуется оценить сходство результатов, полученных из разных экспериментов или алгоритмов кластеризации.
\begin{equation}
    AMI(U, V) = \frac{MI(U, V) - E[MI(U, V)]}{mean(H(U), H(V)) - E[MI(U, V)]}.
\end{equation}

\paragraph{Accuracy} используется для оценки качества классификации, рассматривая степень совпадения найденных и референтных значений. Для оценки качества кластеризации без учителя использование Accuracy может не давать корректных результатов, так как она опирается только на номера меток, которые в кластеризации могут различаться из-за различной нумерации кластеров. Однако при использовании подхода кластеризации с частичным привлечением учителя, где порядок нумерации кластеров совпадает с референтными значениями, данная метрика отражает насколько точно разделяются объекты на кластеры (классы).